\subsection{由两个集合组成的覆盖的特殊情形}

设 X 是道路连通的拓扑空间，B 与 C 是 X 的非空道路连通子集。此外，假设下列两个假设之一成立：

\begin{enumerate}
\def\labelenumi{(\roman{enumi})}
\item
  集合 B 与 C 的内部覆盖 X；
\item
  集合 B 与 C 在 X 中闭，其并等于 X，空间 X、B、C 都是可解圈的，且空间 B ∩ C 局部道路连通。
\end{enumerate}

在这些假设之下，族 (B, C) 的和空间到空间 X 的规范映射满足性质 (VK)（参见 IV, 第~409~页）。

令 A = B ∩ C。由于空间 X 连通，集合 A 非空。设 a 是 A 的一点；令 \(j_{0}\) 为 a 在 A 中的道路连通分支。对 A 的每个异于 \(j_{0}\) 的道路连通分支 j，选取 j 的一点 \(a_{j}\)、B 中一条道路的类 \(\beta_{j}\) 与 C 中一条道路的类 \(\gamma_{j}\)，两者都连接 \(a_{j}\) 与 a；用 \(\varphi_{j}\colon\pi_{1}(A,a_{j})\to\pi_{1}(B,a)\) 与 \(\psi_{j}\colon\pi_{1}(A,a_{j})\to\pi_{1}(C,a)\) 表示由

\[
\varphi_ {j} (v) = \beta_ {j} ^ {- 1} v \beta_ {j} \quad \text {and} \quad \psi_ {j} (v) = \gamma_ {j} ^ {- 1} v \gamma_ {j}
\]

所定义的群同态，其中 \(v \in \pi_1(A, a_j)\)。我们再用 \(\varphi_0\) 与 \(\psi_0\) 表示分别由规范单射导出的、从 \(\pi_1(A, a)\) 到 \(\pi_1(B, a)\) 与 \(\pi_1(C, a)\) 的同态。用 \(\iota_B\) 与 \(\iota_C\) 表示分别由规范单射导出的、从 \(\pi_1(B, a)\) 与 \(\pi_1(C, a)\) 到 \(\pi_1(X, a)\) 中的同态。最后，令 \(\mu\) 为从自由群 \(F(\pi_0(A) - \{j_0\})\) 到 \(\pi_1(X, a)\) 中的唯一同态，使得 \(\mu(j) = \beta_j^{-1}\gamma_j\)（对所有 \(j \in \pi_0(A) - \{j_0\}\)）。

命题 2. --- 存在唯一的群同态

\[
\mathsf {M} \colon \pi_ {1} (\mathrm{B}, a) * \pi_ {1} (\mathrm{C}, a) * \mathrm{F} (\pi_ {0} (\mathrm{A}) - \{j _ {0} \}) \to \pi_ {1} (\mathrm{X}, a)
\]

它与 \(\iota_{\mathrm{B}}\) 在 \(\pi_1(\mathrm{B},a)\) 上、与 \(\iota_{\mathrm{C}}\) 在 \(\pi_1(\mathrm{C},a)\) 上、与 \(\mu\) 在 \(\mathrm{F}(\pi_0(\mathrm{A}) - \{j_0\})\) 上一致。这个同态是满射的，其核是包含元素

\[
\varphi_ {j} (v) j \psi_ {j} (v) ^ {- 1} j ^ {- 1},
\]

（对 \(j\in \pi_0(\mathrm{A}) - \{j_0\}\) 与 \(v\in \pi_1(\mathrm{A},a_j)\)）以及元素

\[
\varphi_ {0} (v) \psi_ {0} (v) ^ {- 1}
\]

（对 \(v \in \pi_1(\mathrm{A}, a)\)）的最小正规子群。

族 (B,C) 所定义的 X 的覆盖的框架 \(\Gamma\) 有两个顶点 \(b\) 与 \(c\)，分别对应两个集合 B 与 C。它的箭集是 \(\pi_0(\mathrm{A})\)；这些箭连接顶点 \(b\) 与 \(c\)。与 \(\Gamma\) 的唯一箭为 \(j_0\) 的子箭图相伴的图是 \(\widetilde{\Gamma}\) 的极大树。命题于是由 IV, 第~412~页命题 1 得出。

例. --- 对 \(n \geqslant 1\)，球面 \(S_{n}\) 是两个闭半球面之并，它们同胚于闭球 \(B_{n-1}\)（TG, VI, 第~12~页），且其交等同于球面 \(S_{n-1}\)。对 \(n \geqslant 2\)，球面 \(S_{n-1}\) 道路连通；由此推出 \(S_{n}\) 的庞加莱群是平凡的（参见 I, 第~127~页，例 3）。

球面 \(S_{0}\) 有两个道路连通分支；这样就重新得到圆 \(S_{1}\) 的庞加莱群同构于一个生成元上的自由群。更精确地说，令 B 与 C 为 \(S_{1}\) 与方程为 \(y \geqslant 0\) 与 \(y \leqslant 0\) 的半平面（在平面 \(R^{2}\) 中）的交。令 \(a = (1,0)\)，\(a' = (-1,0)\)；我们有 \(B \cap C = \{a, a'\}\)；其连通分支为 \(j_{0} = \{a\}\) 与 \(j = \{a'\}\)。令 \(\beta\) 为 C 中道路 \(t \mapsto e^{\pi it}\) 的类；若把 C 等同于 \(R^{2}\)，它在 B 中连接 a 与 \(a'\)。同样，令 \(\gamma\) 为道路 \(t \mapsto e^{-\pi it}\)（在 C 中连接 a 与 \(a'\)）的类。道路 \(\beta\gamma^{-1}\) 是 a 处的圈，由 \(t \mapsto e^{2\pi it}\) 给出。由命题 2，它的类生成群 \(\pi_{1}(\mathbf{S}_{1}, a)\)。

推论 1. --- 同态 \(\mu\) 是单射的。更精确地说，存在一个与 \(\mu\) 相伴的收缩映射，它是群同态。

设 \(\rho\) 为从 \(\pi_1(\mathrm{B},a)*\pi_1(\mathrm{C},a)*\mathrm{F}(\pi_0(\mathrm{A})-\{j_0\})\) 到 \(\mathrm{F}(\pi_0(\mathrm{A}) - \{j_0\})\) 中的唯一同态，它在 \(\pi_1(\mathrm{B},a)\) 与 \(\pi_1(\mathrm{C},a)\) 上诱导平凡同态，在 \(\mathrm{F}(\pi_0(\mathrm{A}) - \{j_0\})\) 上诱导恒等同态。令 N 为同态 M 的核。由命题 2，\(\rho (\mathbf{N})\) 化为单位元。因此存在唯一的同态 r，从 \(\pi_1(\mathrm{X},a)\) 到 \(\mathrm{F}(\pi_0(\mathrm{A}) - \{j_0\})\) 中，使得 \(\rho = r\circ M\)。对所有 \(v\in \mathrm{F}(\pi_0(\mathrm{A}) - \{j_0\})\)，我们有 \(\mathsf{M}(v) = \mu (v)\)，因此 \(r\circ \mu\) 是恒等同态。推论得证。

推论 2. --- 若群 \(\pi_{1}(X,a)\) 是平凡的，则集合 \(A = B \cap C\) 道路连通；若它是交换的，则集合 A 至多有两个道路连通分支。

事实上，若 S 是一个集合，则自由群 F(S) 仅当 S 为空集时是平凡的，仅当 Card S ≤ 1 时是交换的。

推论 3. --- 若群 \(\pi_1(\mathrm{X},a)\) 与 \(\pi_1(\mathrm{A},a)\) 都是平凡的，则群 \(\pi_1(\mathrm{B},a)\) 与 \(\pi_1(\mathrm{C},a)\) 也都是平凡的。

由推论 2，集合 A 道路连通。于是群 \(\pi_1(\mathrm{X},a)\) 同构于群 \(\pi_{1}(B,a)\) 与 \(\pi_{1}(C,a)\) 的自由积。它特别包含同构于群 \(\pi_{1}(B,a)\) 与 \(\pi_{1}(C,a)\) 的子群（A, I, 第~83~页）。因此若 \(\pi_{1}(X,a)\) 平凡，这两个群都是平凡的。

